\section{Implementasi dan Efisiensi}
\label{sec:implementasi-ecc}

Dua persoalan praktis tersisa sebelum ECC dapat dipakai. Yang pertama bersifat
mendasar: kurva eliptik bekerja dengan titik, sedangkan pesan adalah rangkaian
bit. Harus ada cara yang dapat diandalkan untuk mengubah pesan menjadi titik
dan mengembalikannya tanpa kehilangan satu bit pun. Yang kedua bersifat
perancangan: karena ECC menawarkan keamanan yang sama dengan kunci yang jauh
lebih pendek, di mana keunggulan itu paling berharga?

\subsection{Pemetaan Langsung Karakter ke Titik}
\label{subsec:pemetaan-langsung}

Cara paling langsung untuk mengubah pesan menjadi titik adalah memetakan setiap
karakter ke sebuah titik pada kurva. Karena ASCII memuat 256 karakter, kurva
yang dipakai harus memiliki paling sedikit 256 titik agar setiap karakter
mendapat titiknya sendiri. Setelah pemetaan itu disepakati, pesan $M =
\text{'ENCRYPT'}$ --- yang nilai ASCII-nya $69$, $78$, $67$, $82$, $89$, $80$,
dan $84$ --- cukup dipecah menjadi tujuh karakter dan setiap nilainya
dipetakan ke titik yang bersesuaian.

Metode itu sederhana, tetapi kurang aman. Pemetaannya tetap untuk setiap
pemakaian, sehingga dua pesan yang sama selalu menghasilkan titik $P_M$ yang
sama. Akibatnya, penyerang yang tidak mampu memecahkan ECDLP sekalipun masih
dapat mengenali pengulangan: bila dua cipherteks memuat titik kedua yang sama,
ia langsung mengetahui bahwa plainteksnya sama. Kelemahan itu sama dengan
kelemahan buku kode, dan alasan yang sama pula yang membuatnya ditinggalkan.

\subsection{Metode Koblitz}
\label{subsec:koblitz}

Cara yang lebih baik diperkenalkan oleh Neal Koblitz
\citep{koblitz1987}. Gagasannya adalah memanfaatkan sifat residu kuadratik yang
sudah dibahas pada Subbagian~\ref{subsec:medan-berhingga}: untuk nilai $x$ yang
dipilih acak, ruas kanan kurva hanya punya peluang sekitar setengah untuk
menjadi kuadrat modulo $p$. Karena itu penyandi dapat mencoba beberapa nilai
$x$ berurutan sampai menemukan satu yang berhasil. Langkah-langkahnya sebagai
berikut.

\begin{enumerate}
    \item Pilih kurva eliptik $y^2 \equiv x^3 + ax + b \pmod{p}$ yang mengandung
          $N$ buah titik, dengan $N$ yang cukup besar.
    \item Kesepakati pengkodean karakter: angka $0, 1, \dots, 9$ menjadi
          $0, 1, \dots, 9$, dan huruf $A, B, \dots, Z$ menjadi $10, 11, \dots,
          35$.
    \item Kodekan setiap karakter pesan menjadi nilai $m$ di antara $0$ dan
          $35$.
    \item Pilih sebuah bilangan bulat $k$ sebagai parameter basis; nilai ini
          disepakati kedua pihak dan tidak perlu dirahasiakan.
    \item Untuk nilai $m$ tertentu, hitung $x = mk + 1$, lalu sulihkan ke dalam
          persamaan kurva dan tentukan nilai $y$ yang memenuhi.
    \item Bila tidak ada nilai $y$ yang memenuhi, ulangi untuk $x = mk + 2$,
          $x = mk + 3$, dan seterusnya, sampai persamaan itu dapat dipecahkan.
\end{enumerate}

Sebagai contoh, ambil kurva

\[
y^2 \equiv x^3 - x + 188 \pmod{751},
\]

yang memiliki $N = 727$ titik. Karakter yang hendak dikodekan adalah huruf
\texttt{B}, yang menurut langkah 2 bernilai $m = 11$. Ambil parameter basis
$k = 20$. Nilai $x$ pertama yang dicoba adalah

\begin{equation}
x = mk + 1 = (11)(20) + 1 = 221 ,
\label{eq:koblitz-x}
\end{equation}

dan hasilnya diperlihatkan pada Tabel~\ref{tab:koblitz-b}.

\begin{table}[htbp]
    \centering
    \caption{Pencarian titik untuk karakter \texttt{B} ($m = 11$, $k = 20$)
    pada kurva $y^2 \equiv x^3 - x + 188 \pmod{751}$}
    \label{tab:koblitz-b}
    \begin{tabularx}{\linewidth}{@{}c c c
                                   >{\raggedright\arraybackslash}X@{}}
    \toprule
    $x$ & $x^3 - x + 188 \bmod 751$ & \textbf{Kuadrat?} & \textbf{Hasil} \\
    \midrule
    $221$ & $456$ & tidak & tidak ada nilai $y$ yang memenuhi \\
    \addlinespace
    $222$ & $446$ & tidak & tidak ada nilai $y$ yang memenuhi \\
    \addlinespace
    $223$ & $266$ & tidak & tidak ada nilai $y$ yang memenuhi \\
    \addlinespace
    $224$ & $673$ & ya    & $y = 248$ (dan $y = 503$) \\
    \bottomrule
    \end{tabularx}
\end{table}

Jadi karakter \texttt{B} dikodekan menjadi titik $(224, 248)$ pada kurva
eliptik. Pencarian itu hanya memerlukan empat percobaan, sesuai dengan dugaan
peluang: setiap percobaan berpeluang sekitar setengah untuk berhasil,
sehingga rata-rata dibutuhkan dua percobaan, dan peluang memerlukan lebih dari
sepuluh percobaan sangat kecil.

Pada proses \emph{decoding}, penerima memegang titik $(x, y)$ dan harus
memulihkan nilai $m$. Karena penyandian memakai $x = mk + 1$ dan pencariannya
hanya menambah paling banyak $k - 1$, nilai $m$ dapat dihitung dengan

\begin{equation}
m = \left\lfloor \frac{x - 1}{k} \right\rfloor .
\label{eq:koblitz-dekode}
\end{equation}

Untuk titik $(224, 248)$ dengan $k = 20$ diperoleh
\[
m = \left\lfloor \frac{224 - 1}{20} \right\rfloor
  = \lfloor 11{,}15 \rfloor = 11 ,
\]
sehingga pesan semula memang huruf \texttt{B}.

Perhatikan mengapa pembulatan ke bawah selalu memberi jawaban yang benar.
Nilai $x$ yang diperoleh berada di dalam selang $mk + 1$ sampai $mk + k$, sebab
pencarian berhenti sebelum penambahan mencapai $k$. Karena itu
$(x - 1)/k$ selalu berada di antara $m$ dan $m + 1$, dan pembulatan ke bawahnya
tepat sama dengan $m$. Satu catatan penting: agar pemetaan itu tetap satu-satu,
kurva harus memiliki paling sedikit $36k$ titik, sehingga setiap nilai $m$
memperoleh wilayah $x$ sendiri yang tidak tumpang tindih dengan nilai $m$
berikutnya.

\subsection{Efisiensi dan Aplikasi ECC}
\label{subsec:efisiensi-ecc}

Sekarang perhatikan sisi efisiensinya. Untuk menyalurkan kunci AES-128 melalui
algoritma kunci-publik, RSA memerlukan kunci 3072 bit sedangkan ECC cukup 256
bit. Selisih sebesar $12$ kali itu berdampak pada tiga hal sekaligus:
penyimpanan, lebar pita, dan waktu komputasi. Pada perangkat yang menyimpan
ribuan kunci sesi, selisih penyimpanan itu menumpuk; pada jaringan nirkabel
yang lambat, selisih lebar pita terasa langsung; dan pada pemroses kecil,
perkalian titik yang lebih pendek jauh lebih cepat diselesaikan.

Kurva yang dipakai dalam praktik tidak dipilih sembarangan. Sejumlah
parameter baku telah ditetapkan oleh lembaga standardisasi dan diperiksa
keamanannya secara terbuka \citep{fips1865}; Tabel~\ref{tab:kurva-standar}
memuat sebagian di antaranya.

\begin{table}[htbp]
    \centering
    \caption{Beberapa kurva eliptik baku yang dipakai secara luas}
    \label{tab:kurva-standar}
    \begin{tabularx}{\linewidth}{@{}l c
                                   >{\raggedright\arraybackslash}X@{}}
    \toprule
    \textbf{Kurva} & \textbf{Kunci} & \textbf{Pemakaian utama} \\
    \midrule
    secp256k1 & $256$ bit & Mata uang kripto dan rantai blok; parameternya
    dinyatakan dengan cara yang efisien, tetapi titik basisnya berstruktur
    khusus sehingga penanganannya harus hati-hati \\
    \addlinespace
    NIST P-256 & $256$ bit & Sertifikat digital dan TLS; paling banyak
    diperiksa karena paling banyak dipakai \\
    \addlinespace
    Curve25519 & $255$ bit & Pertukaran kunci X25519 dan tanda tangan Ed25519;
    dirancang agar bebas dari pilihan parameter yang membahayakan dan tahan
    terhadap saluran samping \\
    \addlinespace
    NIST P-384 & $384$ bit & Dokumen tingkat keamanan $192$ bit; dipakai bila
    umur data diperkirakan sangat panjang \\
    \bottomrule
    \end{tabularx}
\end{table}

Baris ketiga Tabel~\ref{tab:kurva-standar} layak diperhatikan. Kurva
Curve25519 dirancang oleh Daniel Bernstein \citep{bernstein2006curve25519}
dengan tujuan yang berbeda dari kurva lainnya: bukan untuk kecepatan semata,
melainkan untuk \emph{menghilangkan pilihan}. Pada kurva NIST, pemakai harus
memilih titik basis dari grup yang cukup besar, dan pilihan yang keliru dapat
melemahkan sistem tanpa disadari. Pada Curve25519, titik basisnya ditetapkan
oleh rancangan, dan setiap pemakaian memakai titik yang sama. Sejumlah
pemeriksaan juga tidak diperlukan lagi karena telah terbukti selalu terpenuhi,
sehingga implementasinya lebih sederhana dan lebih kecil peluangnya mengandung
kesalahan.

\subsection{Mengapa ECDLP Lebih Sukar daripada Faktorisasi}
\label{subsec:mengapa-ecdlp-sukar}

Pertanyaan yang wajar muncul pada titik ini adalah mengapa kunci ECC dapat
lebih pendek. Jawabannya terletak pada perbedaan algoritma penyerang terbaik
yang diketahui bagi masing-masing persoalan.

Untuk logaritma diskrit \emph{biasa} --- persoalan yang mendasari
Diffie--Hellman dan ElGamal --- terdapat algoritma \textbf{index calculus}.
Algoritma itu memanfaatkan kenyataan bahwa bilangan bulat dapat difaktorkan
dan disusun kembali, sehingga memungkinkan pencarian yang jauh lebih cerdas
daripada penelusuran menyeluruh \citep{menezes1996}. Versi terbaiknya, yaitu
saringan medan bilangan, berjalan dalam waktu subeksponensial terhadap panjang
kunci. Inilah alasan mengapa Diffie--Hellman dan ElGamal di atas bilangan bulat
memerlukan modulus 3072 bit untuk memperoleh keamanan setara 128 bit.

Algoritma index calculus tidak dapat dipindahkan ke kurva eliptik. Tidak ada
padanan gagasan ``memfaktorkan'' sebuah titik, sebab titik pada kurva bukan
bilangan yang dapat dipecah, melainkan pasangan koordinat yang ditentukan oleh
persamaan kurva. Memang ada usaha memindahkan gagasan itu melalui apa yang
disebut \emph{serangan transfer} --- yaitu memetakan kurva eliptik ke kurva
berorde lebih kecil atau ke medan berhingga biasa --- tetapi serangan semacam
itu hanya berhasil pada kelas kurva yang sempit, dan justru itulah alasan
sejumlah kurva dilarang dipakai. Untuk kurva yang dipilih dengan benar,
penyerang terbaik yang diketahui hanyalah algoritma rho Pollard
\citep{pollard1978}, yang biayanya sebanding dengan akar kuadrat orde grup.
Tabel~\ref{tab:ecdlp-vs-faktorisasi} membandingkan keduanya.

\begin{table}[htbp]
    \centering
    \caption{Biaya penyerangan terbaik yang diketahui pada berbagai panjang
    kunci}
    \label{tab:ecdlp-vs-faktorisasi}
    \begin{tabularx}{\linewidth}{@{}c
                                   >{\raggedright\arraybackslash}X
                                   >{\raggedright\arraybackslash}X@{}}
    \toprule
    \textbf{Keamanan} & \textbf{Kurva eliptik (rho Pollard)} &
    \textbf{RSA (saringan medan bilangan)} \\
    \midrule
    $80$ bit  & $2^{40}$ langkah; kunci $160$ bit  & kunci $1024$ bit \\
    \addlinespace
    $128$ bit & $2^{64}$ langkah; kunci $256$ bit  & kunci $3072$ bit \\
    \addlinespace
    $256$ bit & $2^{128}$ langkah; kunci $512$ bit & kunci $15360$ bit \\
    \bottomrule
    \end{tabularx}
\end{table}

Perhatikan pola pada Tabel~\ref{tab:ecdlp-vs-faktorisasi}. Untuk memperoleh
keamanan $128$ bit pada ECC, penyerang harus menempuh $2^{64}$ langkah, dan
itu berarti panjang kunci cukup dua kali lipatnya, yaitu $256$ bit. Untuk
memperoleh keamanan yang sama pada RSA, penyerang hanya perlu menjalankan
saringan medan bilangan yang biayanya tumbuh jauh lebih lambat terhadap
panjang kunci --- sehingga panjang kunci RSA harus dua belas kali lipat.
Ringkasnya, ECC memperoleh keamanannya ``lebih murah'' karena persoalan yang
mendasarinya lebih merata: tidak ada struktur tersembunyi pada titik-titik
kurva yang dapat dimanfaatkan penyerang.

Yang perlu dicatat adalah bahwa keunggulan itu hanya berlaku bagi kurva yang
dipilih dengan benar. Kurva yang orderya dapat dihitung dengan mudah, kurva
dengan orde yang habis dibagi bilangan kecil, dan kurva yang dapat ditransfer
ke medan berhingga biasa semuanya runtuh jauh lebih cepat daripada dugaan
umum. Karena itu perkakas verifikasi pada bagian praktikum di akhir bab ini
tidak hanya menjalankan aritmetika kurva, tetapi juga memeriksa orde titik dan
menjalankan penyerangan logaritma diskrit berbiaya akar orde pada kurva mainan
--- untuk memperlihatkan kepada pembaca seberapa cepat keamanan itu hilang
ketika orderya menyusut.
